\chapter{Introduction}

\section{Background}

\subsection{What Is Stimulated Brillouin Scattering?}

Stimulated Brillouin Scattering (SBS) is a nonlinear optical effect that occurs due to the interaction between light and sound within a material 
\parencite{eggletonBrillouinScatteringPart2022}. When a strong optical (pump) wave, typically from a laser, travels through a medium, 
it undergoes Brillouin scattering as it interacts with co-propagating acoustic (density) waves generated within the medium through 
changes in density induced by strong electric fields. These acoustic waves form a moving refractive index grating 
that scatters part of the light backwards, producing a backward-propagating, frequency-shifted Stokes wave. 
The interaction between these three waves creates a strong self-reinforcing amplification of the Stokes and acoustic waves, 
which leads to exponential growth of the Stokes wave as energy is progressively drained from the pump.

The Brillouin scattering effect most commonly occurs due to the acoustic waves being optically excited. 
The two main physical effects related to this scattering are as follows: first, acoustic waves compress the medium, which alters 
its optical permittivity via photoelasticity. Second, acoustic waves deform the physical boundaries of the medium, which directly 
alters the permittivity at its edges and interfaces.

\begin{figure}
  \centering
  \thesisimage[width=0.85\textwidth]{SBS-picture-with-the-levels}
  \caption{Pump and Stokes wave interaction through time \parencite{baiStimulatedBrillouinScattering2018}.}
  \label{fig:SBSWithLevels}
\end{figure}

A key result of this scattering effect is that the backward-propagating Stokes wave is frequency-shifted relative to the pump. 
This is known as the Brillouin Frequency Shift (BFS). It arises because the acoustic waves (acting as an acoustic grating) are moving, 
so the scattered light experiences a Doppler shift and momentum change. In silica fibre at a wavelength of \qty{1550}{\nano\metre}, 
this downshift is typically around \qty{11}{\giga\hertz} \parencite{agrawalNonlinearFiberOptics2019}. The pump, the acoustic wave and 
the Stokes wave therefore form a feedback loop, and it is this self-reinforcing amplification that the \enquote{stimulated} in SBS refers to.

\subsection{Applications of SBS}

SBS has many important applications in modern photonics and can act as both a useful tool and a limiting constraint in systems 
\parencite{baiStimulatedBrillouinScattering2018}.

In optical signal processing, SBS can help create narrow-linewidth fibre lasers and amplifiers, which are essential for 
high-precision sensing and signal processing in systems where narrow linewidths are important for signal coherence and 
minimising noise \parencite{agrawalNonlinearFiberOptics2019}. Another important application is detecting changes in the 
environment such as temperature and pressure. This capability makes SBS useful for distributed sensing, such as monitoring the 
structural health of civil infrastructure. The main approach for such sensing is known as Brillouin Optical Time Domain Analysis.

Brillouin scattering can also have a negative effect on optical systems. 
Backscattering can occur due to SBS in high-power fibre optic systems, limiting the peak power that can be transmitted without 
distortion or noise. These effects increase with propagation distance, leading to further distortion. Furthermore, in high-power 
fibre lasers and amplifiers SBS acts as the dominant effect that caps the output power. This happens once the pump exceeds the SBS 
threshold and energy is redirected into the backward-propagating Stokes wave rather than the intended output.

Whether trying to suppress or use SBS, accurate simulation of its dynamics is important for understanding the effects 
within a system. These effects rarely have analytic solutions, thus motivating the need for a generalised simulation framework.

\subsection{The Main Equations of SBS}

SBS dynamics are governed by the three-wave interaction equations, \cref{eq:sbs}:
\begin{subequations}
\label{eq:sbs}
\begin{alignat}{2}
  \left(\frac{\partial}{\partial T} + V_1\frac{\partial}{\partial Z} + \mu_1\right) a_1 &= -\kappa\, a_2 b,  &\qquad& \text{(Pump)} \label{eq:pump}\\
  \left(\frac{\partial}{\partial T} + V_2\frac{\partial}{\partial Z} + \mu_2\right) a_2 &= \kappa\, a_1 b^*, && \text{(Stokes)} \label{eq:stokes}\\
  \left(\frac{\partial}{\partial T} + V_3\frac{\partial}{\partial Z} + \mu_3\right) b   &= \kappa\, a_1 a_2^*. && \text{(Acoustic)} \label{eq:acoustic}
\end{alignat}
\end{subequations}

These equations are a system of coupled, nonlinear partial differential equations, which describe the evolution of the three main 
wave envelopes. 
Here $a_1$, $a_2$ and $b$ are the slowly varying envelopes of the pump, Stokes and acoustic fields respectively. 
They are functions of position $Z$ along the waveguide and time $T$. $V_1$, $V_2$ and $V_3$ are the velocities at which each envelope 
travels along the waveguide, $\mu_1$, $\mu_2$ and $\mu_3$ are the corresponding linear damping rates, and $\kappa$ is the 
coupling coefficient that sets the strength of the interaction between the waves. The asterisk denotes the complex conjugate. 

Each equation shares the same structure, where the left-hand side contains a time derivative, a spatial derivative scaled by the 
respective wave's velocity, and a linear damping term. The right-hand side captures the nonlinear coupling between the three waves through the coupling coefficient $\kappa$. Energy from the pump is drained by the Stokes 
and acoustic waves, hence the right-hand side of \cref{eq:pump} carries a negative sign. These equations represent a simplified three-wave 
model that neglects the optical Kerr effect and acoustic noise, to isolate the core dynamics of SBS 
\parencite{montesThreeWaveDissipativeBrillouin2008}.

\Cref{fig:PulseSimulated} illustrates the kind of result sought: given the envelopes at the start, the simulation evolves them 
to give their values at later times.

\begin{figure}
  \centering
  \thesisimage[width=0.85\textwidth]{SimpleGaussianProp}
  \caption{A Gaussian pulse before (left) and after (right) linear propagation. Solid lines show the envelope amplitude 
and dashed lines show its phase. The simulation takes the initial state as input and computes the evolved state.}
  \label{fig:PulseSimulated}
\end{figure}

Numerical simulation of SBS dynamics can be set up in two ways: following how the waves change in space as time passes, or following 
how they change in time as they travel along the waveguide. The equations here use the first. They describe how the shape of each wave 
along the waveguide changes through time. This is the more natural choice for the short structures considered here, where light 
enters and leaves through the ends, because the state of the whole waveguide can be described at one instant and stepped forward in 
time (\cref{fig:PulseSimulated}).

\section{Why Simulation Is Required}

Accurately modelling SBS requires solving a set of coupled nonlinear partial differential equations which define the characteristics 
of the waves. Analytical solutions exist only for simplified cases, such as steady-state operation, where the waves do not change over 
time \parencite{agrawalNonlinearFiberOptics2019}. Once the pump is pulsed or additional effects such as the optical Kerr nonlinearity 
are included, the nonlinear coupling terms in \cref{eq:pump,eq:stokes,eq:acoustic} prevent the equations from being decoupled, and no 
general closed-form solution is known. These cases therefore require numerical simulation.

Numerical simulation overcomes these limitations by providing access to the entire evolution of all three wave envelopes, 
including quantities that are difficult or impossible to measure inside a physical waveguide. It also allows parameters to be varied 
freely and cheaply, enabling researchers to explore system designs, test suppression strategies or sensing configurations, and compare 
theory against experiment before committing to costly fabrication or measurement. The closed-form solutions that do exist for specific 
setups also provide a practical means of testing and validating the simulation's numerical accuracy.


\section{The Existing C++ Codebase}

An existing C++ codebase for the numerical simulation of SBS, created in November 2025 by Professor Michael Steel, provides the basis
on which this research builds. This existing codebase builds and runs on Linux, Windows, and macOS, and is written in C++23 for the core
of the simulation, along with some Python scripts for testing and setup.
The core currently consists of about 14,000 lines of C++ across 97 source files, excluding third-party code, and builds on Eigen for linear algebra, FFTW (MKL on Windows) for Fourier transforms, toml++ for reading the control file, and Catch2 for testing.
It simulates how light pulses travel through an optical waveguide and accounts for SBS and Kerr-nonlinear effects.
Plots of these pulses are output at user-defined intervals of the simulation for analysis, and it can simulate many waves, both forward-propagating
and backward-propagating at the same time.

\begin{figure}
  \centering
  \thesisimage[width=0.85\textwidth]{codebase-architecture}
  \caption{Current codebase architecture schematic.}
  \label{fig:codebase-architecture}
\end{figure}


\Cref{fig:codebase-architecture} shows a very high-level overview of the codebase architecture.
A TOML file configures the specifics of the situation. This file can be used to change a variety of aspects of the simulation, including which linear and nonlinear effects are active, the shape of each forward- and backward-propagating pulse (one of three preset shapes), and the analysis settings.
This file is then parsed internally to set up the starting conditions for the simulation in the Parameters and Simulation Context stages.
The simulation supports a user choice of propagation methods which define how the effects are applied to the pulses at each time step.
The chosen solver (propagation method) is used in the Propagation stage for every step of the simulation, and defaults to symmetric split-step if omitted from the configuration file.
Throughout propagation, the pulse may periodically be printed on a separate thread for analysis, then once the propagation finishes, the final analysis is done in the Analysis stage.

Most components in the Parameters and Simulation Context stages are default-constructed and then configured by a separate
\texttt{setup()} call once the parameters they depend on have been parsed.
This two-phase initialisation departs from standard C++ practice, under which
a constructor should leave an object fully initialised and usable
\parencite{CoreGuidelines}. Here, an object is in an unusable state until
\texttt{setup()} has been called. The pattern is nonetheless common in
scientific computing libraries such as PETSc, where objects are created,
configured, and then explicitly set up in separate steps \parencite{UserGuidePETSc3261}.
It allows components to be declared as plain members of the classes that own
them and configured later, in the order their dependencies require.

Physical models of SBS involve quantities spanning many orders of magnitude, such as $10^{-27}\,\mathrm{s^2/m}$, alongside lengths in metres.
Entered directly into a computer, these values risk overflow and underflow, and make it hard to tell whether the behaviour of the simulation is correct.
These differences in scale are only a side effect of the chosen units: rescaling each quantity by a characteristic value, known as non-dimensionalisation, brings them back to similar sizes.
Theorists commonly work this way, while experimentalists generally prefer to work with physical units.
The codebase supports both through the \texttt{solver.normalise\_quantities} option in the TOML file.

\section{Propagation Methods}

The existing SBS simulation has three different propagation methods: symmetric split-step, interaction-picture fourth-order Runge--Kutta (RK4), 
and interaction-picture Cash--Karp Runge--Kutta (RK5). All three methods solve the same coupled slowly varying envelope equations but differ 
in how they apply the linear and nonlinear effects.
Writing the linear and nonlinear operators as $\mathcal{D}$ and $\mathcal{N}$, advancing the field over a step $h$ is equivalent to
applying $e^{h(\mathcal{D}+\mathcal{N})}$. The operators stay the same throughout the simulation, and only the field is updated at each step. On the numerical grid these operators become matrices that do not commute, so
$e^{h\mathcal{D}}e^{h\mathcal{N}} \neq e^{h(\mathcal{D}+\mathcal{N})}$, but the two become very close when $h$ is small, which is why the
methods take many short steps.

The symmetric split-step method advances the field from $z_k$ to $z_{k+1}$ by applying a half step ($h_k/2$) 
of linear effects, followed by a full step ($h_k$) of nonlinear effects starting from that result, 
then a second half step of linear effects. \Cref{fig:symmetric-split-step} illustrates the symmetric splitting: 
the red $\mathcal{D}$ arrows represent linear half steps, while the blue $\mathcal{N}$ arrows represent full nonlinear steps, 
and the dashed arrows carry each result to the start of the next operator.
Although the nonlinear step is evaluated using RK4, the overall symmetric split-step method is second-order accurate in time.

\begin{figure}
  \centering
  \thesisimage[width=0.8\textwidth,trim=0 0 30 0,clip]{symmetric-split-step}
  \caption[Symmetric split-step propagation]{Symmetric split-step propagation, with computational step $k$ shown in the dashed box \parencite{balacMathematicalAnalysisAdaptive2014}.}
  \label{fig:symmetric-split-step}
\end{figure}

The RK4 method is fourth-order accurate in time, so halving the time step is expected to reduce the global temporal error by approximately 
a factor of 16. In comparison, the second-order symmetric split-step method reduces its error by approximately a factor of four. 
However, RK4 requires more function evaluations per step, meaning its higher convergence order does not necessarily translate to greater 
computational efficiency.

\begin{equation}
  \label{eq:rk4}
  \begin{aligned}
    A^{(i)} &= U_{h/2}^{\mathcal{D}} A_n, \\
    k_1 &= U_{h/2}^{\mathcal{D}} f(A_n, z_n), \\
    k_2 &= f\!\left(A^{(i)} + \tfrac{h}{2}k_1,\, z_n + \tfrac{h}{2}\right), \\
    k_3 &= f\!\left(A^{(i)} + \tfrac{h}{2}k_2,\, z_n + \tfrac{h}{2}\right), \\
    k_4 &= f\!\left(U_{h/2}^{\mathcal{D}}\left(A^{(i)} + h k_3\right),\, z_n + h\right), \\
    A_{n+1} &= U_{h/2}^{\mathcal{D}}\left[A^{(i)} + \tfrac{h}{6}\left(k_1 + 2k_2 + 2k_3\right)\right] + \tfrac{h}{6}k_4.
  \end{aligned}
\end{equation}

The fourth-order Runge--Kutta method evaluates the nonlinear derivatives at four intermediate stages, which are then combined using a weighted average and applied to the field, as shown in \cref{eq:rk4}. The linear part of the problem is easy: because $\mathcal{D}$ is diagonal in Fourier space, dispersion alone can be solved arbitrarily far ahead in a single step. The nonlinear part cannot be diagonalised and has no such solution except in special cases, so the whole problem cannot be solved in one step. The interaction picture takes advantage of this by instead moving the easy linear part onto the operators, the $U^{\mathcal{D}}_{h} = e^{h\mathcal{D}}$ terms in \cref{eq:rk4}, and leaving only the hard nonlinear part on the field for RK4 to integrate, which reduces the work done per step. Unlike the symmetric split-step method, the RK4 implementation incorporates the linear effects throughout, allowing the interaction between linear and nonlinear effects to be accounted for within each step. This requires four nonlinear derivative evaluations per step, increasing the computational cost per step compared to simpler integration schemes.

Similarly, the RK5 method of propagation is fifth-order accurate in time, so halving the time step reduces the temporal error by approximately a factor of 32. This comes at the cost of additional calculations at each step. \Cref{eq:rk5-stages} in \cref{app:rk5} gives all six stages of the RK5 method, and \cref{eq:rk5-solution} the fifth-order solution built from them. The Cash--Karp scheme also embeds a fourth-order solution built from the same six stages, and the difference between the two provides a local error estimate that can be used to adapt the step size \parencite{balacMathematicalAnalysisAdaptive2014}.


\section{The Research Gap}

The focus of this research project is to improve the simulation software in terms of three main criteria: performance, modularity and portability. Program performance is fundamental for timely running of large simulations and reducing the time spent waiting for results. Improving the modularity of the program means it can be easily extended for further research without compromise, for example by adding a new nonlinear effect in one localised part of the code. Portability is important for the program to not just be able to run on any device but to be able to do so efficiently.

Current approaches to simulation software typically focus on creating instruments useful for either theoretical physicists or experimental physicists, rarely addressing the needs of both. Theoretical research focuses on working with idealised or unattainable conditions, whereas experimental research is constrained by real-world limitations. This research project will address the need to create software that meets the requirements of both sides without compromising on either.

When well-developed libraries for simulation do not exist, the common practice in research is to write ad-hoc Python scripts tailored to a specific experimental configuration, use them once, and discard or abandon them for the next task. While Python allows for rapid prototyping, these scripts accumulate technical debt quickly and consequently are rarely used beyond their original purpose. Python's interpreted nature also restricts the performance of any simulations run. This means that each new research task effectively starts from scratch, with researchers reinventing core SBS simulation logic repeatedly rather than building on a stable foundation.

\section{Research Significance and Physics Outcomes}

The significance of this research lies in producing an enhanced simulation codebase for SBS that can serve as a reusable and extensible foundation for future research. Different research requires different environmental configurations, such as pulse shapes and other nonlinear and linear effects. A flexible and tunable library allows researchers to easily adapt the software to the specific needs of their research, cutting out the time that would be spent rewriting core SBS logic for each specific research task. Well-written and flexible code also means that as new discoveries occur, they can be added to the library without large architectural changes.

From a performance standpoint, the aim of this research project is for all simulations to run faster, leading to less dead time for researchers whilst waiting for results. This outcome is most important for simulations done at a large scale. By writing well-optimised C++ and GPU (graphics processing unit) accelerated code, not only does the code become more efficient for all simulations, but it can also make it practical to simulate conditions that were previously too large to simulate, potentially revealing unseen effects.

To increase the modularity of the system, on top of writing clean C++ that follows best coding practices, the research will add custom function definitions that can be used to define waves and other spatially dependent fields. This will free researchers from a preset set of options and broaden what can be simulated. Clean coding practices also make the library easy to extend for specific research needs.

To address portability, all changes will consider the ability of the codebase to be developed and executed across different 
operating systems (typically macOS or Linux for theorists and Windows for experimentalists).
This will lead to fewer barriers to using and extending the library for all types of physicists.

\section{Research Objectives}

This research is guided by one central objective: to improve the existing SBS modelling codebase into a library that 
balances performance, modularity for further research, and portability. Success against all three metrics will be 
evaluated through benchmarking, portability testing, and architectural and structural code analysis against software 
engineering best practices.

Performance success is defined as a measurable reduction in wall-clock time for the simulation relative to the 
existing baseline implementation, verified through repeated benchmarking under controlled conditions. 
Modularity success is defined as the ability to add new physics logic in a localised code fashion, 
without modifying unrelated parts of the code, verified through architectural review and testing. Portability success is 
defined as the ability to build and run the codebase across macOS, Windows and Linux without any code changes, 
producing outputs that agree to within floating-point tolerance.

\section{CPU and GPU Computing}

A CPU (central processing unit) is a hardware component responsible for executing computing instructions of a 
specific instruction set architecture (ISA), performing computations and modifying architectural state. 
Modern CPUs contain multiple cores, each of which has its own levels of local cache and is capable of independently 
executing instructions. CPUs are designed for high single-core performance and low execution latency. 
To achieve this, they use a wide range of techniques such as SIMD (single instruction, multiple data) instructions, pipelined instruction processing, 
branch prediction, out-of-order execution, and multi-level cache hierarchies \parencite{stokesMachine2006}.

A GPU, however, is a more specialised piece of hardware originally designed for rendering graphics, 
which now sees widespread use in machine learning, high-performance computing and other computationally intensive workloads. 
GPUs are designed for highly parallel computation, containing thousands of arithmetic execution units organised into multiple 
processing groups. Rather than maximising the computational speed of a single thread, GPUs are designed to maximise 
computational throughput by executing large numbers of threads simultaneously. This makes a GPU well suited to large, 
data-parallel operations where the same computation is applied across many elements \parencite{nickollsGPUComputingEra2010}.

A GPU operates on data held in its own local memory, known as Video Random Access Memory (VRAM), which is 
physically separate from the system RAM used by the CPU. Data produced on one device and needed by the other must 
therefore be copied between the two memories over the interconnect that links them, typically a PCIe bus. 
As illustrated in \cref{fig:gpu-cpu-architecture}, each transfer introduces latency on the order of microseconds, 
and the available bandwidth is limited to approximately \qty{25}{\giga\byte\per\second} for a PCIe~4.0~x16 connection, 
considerably lower than the bandwidth of the GPU's own memory. The total transfer time therefore depends on both 
the fixed overhead per transfer and the amount of data being transferred, so workloads that frequently move data 
between the CPU and GPU can lose much of the benefit of GPU acceleration.

\begin{figure}
  \centering
  \thesisimage[width=0.75\textwidth]{gpu-cpu-architecture}
  \caption{High-level GPU and CPU architecture.}
  \label{fig:gpu-cpu-architecture}
\end{figure}
